%%%%%%%%%%  scan idx 25  |  folio 12  %%%%%%%%%%
essentiellement non recollables. Voici une construction g\'en\'erale de $L^{X/Y}_\bullet$, valable plus g\'en\'eralement pour tout morphisme $f:X\to Y$ de topos (commutativement) annel\'es~: Soit $B=\underline O_X$, $A=f^{-1}(\underline O_Y)$, de sorte que $A$ est un Anneau sur l'espace $X$, et $B$ une $A$-Alg\`ebre. Soit $T\to B$ un homomorphisme de faisceaux d'\underline{ensembles} qui engendre $B$ comme $A$-Alg\`ebre, i.e.\ tel que l'homomorphisme de $A$-Alg\`ebres
\[
\phi:C=A[T]\longrightarrow B
\]
correspondant soit un \'epimorphisme de faisceaux d'ensembles, o\`u $A[T]$ d\'esigne la $A$-Alg\`ebre commutative libre engendr\'ee par $T$ (i.e.\ le faisceau associ\'e au pr\'efaisceau $U\longmapsto A(U)[T(U)]$ = $A(U)$-alg\`ebre commutative libre engendr\'e par $T(U)$).
Soit $J$ le noyau de $\phi$, et consid\'erons comme dans 4.3 le complexe de chaines de longueur $1$ d\'efini par l'homomorphisme naturel
\[
J/J^2\xrightarrow{d}\Omega^1_{C/A}\otimes_C B.
\]
Il n'est pas difficile de voir que, \`a isomorphisme canonique pr\`es dans la cat\'egorie $D(X)$, le complexe ainsi construit est ind\'ependant du choix de $T$ et de $T\to B$, de sorte qu'il m\'erite la d\'esignation de \uline{complexe cotangent relatif $L^{X/Y}_\bullet$ de $X$ sur $Y$}. On v\'erifie \'egalement que pour tout ouvert affine de $X$ au dessus d'un ouvert affine de $Y$, il est donn\'e (\`a isomorphisme canonique pr\`es) par le proc\'ed\'e de 4.3.

Dans le cas o\`u $f$ est localement d'intersection compl\`ete, on voit de plus que l'on obtient m\^eme un complexe \underline{parfait}, ce qui permet alors de d\'efinir encore un \'el\'ement $T_f$ dans $K^\bullet(X)$ par la formule (4.3). Signalons qu'ici encore, il n'aurait pas \'et\'e possible en g\'en\'eral de d\'efinir $T_f$ comme un \'el\'ement de $K^\bullet_{\mathrm{naif}}(X)$, et rien ne permet de supposer que $T_f$ soit \underline{globalement} isomorphe dans $D(X)$ \`a un complexe born\'e \`a composantes localement libres de type fini. Cel\`a est donc une justification suppl\'ementaire du point de vue avanc\'e dans 4.1 en faveur du $K^\bullet(X)$ «~sophistiqu\'e~», de pr\'ef\'erence \`a $K^\bullet_{\mathrm{naif}}(X)$.

\subsubsection*{4.5.}
Sous r\'eserve de la question soulev\'ee dans 4.2, nous avons donc encore r\'euni tous les \'el\'ements de structure pour donner un sens \`a (1.2) pour un morphisme propre localement d'intersection compl\`ete de sch\'emas quelconques (l'hypoth\`ese noeth\'erienne faite implicitement dans les sections qui pr\'ec\`edent pouvant en fait s'\'eliminer sans difficult\'e essentielle). Cela ne signifie pas malheureusement que, m\^eme sous cette r\'eserve, la formule en question soit \'etablie, m\^eme dans le cas o\`u $Y$ est le spectre d'un corps 
